\documentclass[../../main.tex]{subfiles}
\graphicspath{{\subfix{../../image/}}}

\begin{document}

\section{偏微分方程}\label{section:ch01-pde}

本章综述求解偏微分方程所涉及的主要理论问题。

为阅读后续讨论，读者最好先翻阅\cref{chapter:appA-记号} 中给出的记号，特别是偏导数的多重指标记号。

一个偏微分方程（PDE）是含有一个未知的多元函数及其某些偏导数的方程。利用\cref{chapter:appA-记号} 中说明的记号，我们可以符号地写出一个典型的 PDE 如下。取定整数 $k\geqslant 1$，设 $U$ 为 $\mathbb{R}^n$ 的一个开子集。

\begin{definition}\label{definition:ch01-k阶pde}
形如
\begin{align}\label{align:ch01-pde-general}
F(D^k u(x), D^{k-1}u(x), \cdots, Du(x), u(x), x)=0 \qquad (x\in U),
\end{align}
的表达式称为一个 $k$ 阶偏微分方程，其中
\begin{align*}
F:\mathbb{R}^{n^k}\times\mathbb{R}^{n^{k-1}}\times\cdots\times\mathbb{R}^n\times\mathbb{R}\times U\to\mathbb{R}
\end{align*}
是给定的，而
\begin{align*}
u:U\to\mathbb{R}
\end{align*}
是未知函数。
\end{definition}

所谓求解该 PDE，就是要找出所有满足\eqref{align:ch01-pde-general} 的 $u$，这或许只是在那些于 $\partial U$ 的某一部分 $\Gamma$ 上满足某些辅助边界条件的函数中去找。所谓求解，理想地说，就是得到简单、显式的解；若做不到这一点，则至少要推导出解的存在性及其他性质。

\begin{definition}\label{definition:ch01-方程分类}
\begin{enumerate}[(i)]
\item 若偏微分方程\eqref{align:ch01-pde-general} 具有形式
\begin{align*}
\sum_{|\alpha|\leqslant k} a_\alpha(x) D^\alpha u = f(x),
\end{align*}
其中 $a_\alpha$（$|\alpha|\leqslant k$）与 $f$ 是给定函数，则称它为\textbf{线性}的。若 $f\equiv 0$，则称此线性 PDE 是齐次的。
\item 若 PDE \eqref{align:ch01-pde-general} 具有形式
\begin{align*}
\sum_{|\alpha|=k} a_\alpha(x) D^\alpha u + a_0(D^{k-1}u, \cdots, Du, u, x)=0,
\end{align*}
则称它为\textbf{半线性}的。
\item 若 PDE \eqref{align:ch01-pde-general} 具有形式
\begin{align*}
\sum_{|\alpha|=k} a_\alpha(D^{k-1}u, \cdots, Du, u, x) D^\alpha u + a_0(D^{k-1}u, \cdots, Du, u, x)=0,
\end{align*}
则称它为\textbf{拟线性}的。
\item 若 PDE \eqref{align:ch01-pde-general} 非线性地依赖于最高阶导数，则称它为\textbf{完全非线性}的。
\end{enumerate}
\end{definition}

通俗地说，偏微分方程组就是若干个未知多元函数所对应的若干个 PDE 的集合。

\begin{definition}\label{definition:ch01-方程组}
设 $k$ 为正整数。形如
\begin{align}\label{align:ch01-system-general}
\mathbf{F}(D^k\mathbf{u}(x), D^{k-1}\mathbf{u}(x), \cdots, D\mathbf{u}(x), \mathbf{u}(x), x)=0 \qquad (x\in U),
\end{align}
的表达式称为一个 $k$ 阶偏微分方程组，其中
\begin{align*}
\mathbf{F}:\mathbb{R}^{mn^k}\times\mathbb{R}^{mn^{k-1}}\times\cdots\times\mathbb{R}^{mn}\times\mathbb{R}^m\times U\to\mathbb{R}^m
\end{align*}
是给定的，而
\begin{align*}
\mathbf{u}:U\to\mathbb{R}^m,\qquad \mathbf{u}=(u^1,\cdots,u^m)
\end{align*}
是未知函数。
\end{definition}

这里我们假定方程组中数量与未知函数 $(u^1,\cdots,u^m)$ 相同，即含有同样数目的 $m$ 个标量方程。这是最常见的情形，不过其他方程组的方程个数也可能少于或多于未知函数个数。方程组按显而易见的方式分类为线性、半线性等。

\textbf{记号.} 我们用``PDE''同时作为单数``偏微分方程''与复数``偏微分方程组''的缩写。













\end{document}